% !TEX program = xelatex
% 공통수학2 · Ⅰ. 도형의 방정식 · 01 선분의 내분 · 1/3차시
% 학생용 활동지 (답안 없음)
% 컴파일: XeLaTeX (Overleaf 메뉴에서 Menu > Compiler > XeLaTeX 로 지정)

\documentclass[11pt,a4paper]{article}

\usepackage{kotex}
\usepackage{iftex}
\ifPDFTeX\else
  \setmainhangulfont{Noto Serif CJK KR}
  \setsanshangulfont{Noto Sans CJK KR}
\fi

\usepackage[margin=2.1cm]{geometry}
\usepackage{parskip}
\usepackage{enumitem}
\usepackage{array}
\usepackage{tabularx}
\usepackage{xcolor}
\usepackage{amssymb}
\usepackage{tikz}
\usepackage[most]{tcolorbox}
\usepackage{fancyhdr}
\usepackage{titlesec}

\definecolor{goldc}{HTML}{B07C22}
\definecolor{goldstrong}{HTML}{8A5F16}
\definecolor{indigoc}{HTML}{2B3B57}
\definecolor{indigosoft}{HTML}{6E82A6}
\definecolor{linec}{HTML}{D6DACB}
\definecolor{surface2}{HTML}{E4E8DC}
\definecolor{writeline}{HTML}{C7CCBA}
\definecolor{inksoft}{HTML}{52604F}

\titleformat{\section}
  {\normalfont\sffamily\bfseries\large\color{indigoc}}
  {\thesection}{0.6em}{}
\titlespacing{\section}{0pt}{14pt}{6pt}

% ---------- 재사용 박스 ----------
\newtcolorbox{goalbox}{
  enhanced, breakable, colback=white, colframe=linec, boxrule=0.5pt,
  leftrule=3pt, colframe=goldc, arc=2pt,
  left=10pt, right=10pt, top=8pt, bottom=8pt
}
\newtcolorbox{sheetbox}{
  enhanced, breakable, colback=white, colframe=linec, boxrule=0.6pt,
  arc=3pt, left=11pt, right=11pt, top=9pt, bottom=9pt
}

% ---------- 줄쳐진 기록 공간 (일반 텍스트 흐름 -- 페이지 넘김에 안전, tikz 그림 아님) ----------
% #1 = 줄 개수
\newcommand{\writespace}[1]{%
  \par\nobreak\vspace{3pt}
  \foreach \n in {1,...,#1} {%
    \noindent\textcolor{writeline}{\rule{\linewidth}{0.4pt}}\par\vspace{7pt}%
  }
}

\newcommand{\fillblank}[1]{\makebox[#1]{\hrulefill}}

\pagestyle{fancy}
\fancyhf{}
\renewcommand{\headrulewidth}{0pt}
\fancyfoot[L]{\footnotesize\color{inksoft} 공통수학2 · 2026학년도 2학기 · 지평선고등학교}
\fancyfoot[R]{\footnotesize\color{inksoft} 다음 시간 --- 02 좌표평면 위의 선분의 내분}

\begin{document}

% ---------- 머리 : 이름/반/번호 + 골드 밑줄 ----------
\noindent
\begin{minipage}[b]{0.62\linewidth}
  {\sffamily\footnotesize\color{indigosoft} 공통수학2 · Ⅰ. 도형의 방정식}\\[2pt]
  {\sffamily\bfseries\Large 01 선분의 내분 \textperiodcentered\ 1차시}
\end{minipage}\hfill
\begin{minipage}[b]{0.36\linewidth}
  \raggedleft\small\color{inksoft}
  반\ \fillblank{1.6cm} \quad 번호\ \fillblank{1.6cm}\\[4pt]
  이름\ \fillblank{3.6cm}
\end{minipage}

\vspace{4pt}
\noindent{\color{goldc}\rule{\linewidth}{2.2pt}}
\vspace{10pt}

\begin{goalbox}
\textbf{오늘의 목표}\\
수직선 위에서 `내분'이 무엇인지 이해하고, 두 점을 $m:n$으로 나누는 점의 좌표를 직접 구할 수 있다.
\vspace{4pt}

\small\color{inksoft}
$\square$ 자신 있음 \quad $\square$ 조금 헷갈림 \quad $\square$ 처음 봄 \hfill (수업 전 나의 상태에 표시)
\end{goalbox}

\section{생각 열기 --- 일직선 도로 위의 세 지점}

\begin{sheetbox}
아래는 창원대로(우리나라에서 가장 긴 직선 도로)의 일부를 한 눈금 1\,km로 하여 수직선 위에 나타낸 것이다.
A는 버스 터미널, B는 올림픽 공원, C는 성산구청 근처의 한 지점이다.

\begin{center}
\begin{tikzpicture}[scale=0.78, every node/.style={font=\small}]
  % 실제 값에서 계산한 눈금: 0~9, 10개
  \draw[thick, indigosoft] (0,0) -- (9,0);
  \foreach \x in {0,...,9} {
    \draw[inksoft] (\x,-0.14) -- (\x,0.14);
    \node[below=1pt, font=\tiny, inksoft] at (\x,-0.16) {\x};
  }
  \filldraw[goldc] (1,0) circle (3.6pt);
  \node[above=4pt] at (1,0) {\textbf{A}};
  \node[below=13pt, font=\footnotesize] at (1,-0.16) {터미널};
  \filldraw[goldc] (3,0) circle (3.6pt);
  \node[above=4pt] at (3,0) {\textbf{B}};
  \node[below=13pt, font=\footnotesize] at (3,-0.16) {올림픽공원};
  \filldraw[goldc] (8,0) circle (3.6pt);
  \node[above=4pt] at (8,0) {\textbf{C}};
  \node[below=13pt, font=\footnotesize] at (8,-0.16) {성산구청};
\end{tikzpicture}\\
{\footnotesize\color{inksoft} 한 눈금 = 1\,km \quad\textperiodcentered\quad A는 1번 눈금, B는 3번 눈금, C는 8번 눈금 위에 있다}
\end{center}

\begin{enumerate}[leftmargin=1.4em, itemsep=2pt]
  \item 두 점 A, B 사이의 거리와 두 점 B, C 사이의 거리를 각각 구해 보자.
  \item $\overline{AB}:\overline{BC}$를 가장 간단한 자연수의 비로 나타내 보자.
\end{enumerate}
\writespace{3}

\vspace{6pt}
\textbf{확장 질문} --- 우리가 사는 지평선 지역(김제·만경 일대)에도 이렇게 곧게 뻗은 길이나 둑길이 있을까? 있다면 어떤 길인지, 실제로 어떻게 확인할 수 있을지 적어 보자.
\writespace{2}
\end{sheetbox}

\section{개념 정리 --- 내분이란?}

\begin{sheetbox}
선분 AB 위의 점 P에 대하여
\[
\overline{AP} : \overline{PB} = m : n \quad (m>0,\ n>0)
\]
일 때, 점 P는 선분 AB를 \fillblank{1.2cm} : \fillblank{1.2cm}으로 \fillblank{2.4cm}한다고 하며,
점 P를 선분 AB의 \fillblank{2.4cm}이라고 한다.

\vspace{8pt}
\textbf{함께 채우기} --- 수직선 위의 두 점 A($x_1$), B($x_2$) ($x_1<x_2$)에서 선분 AB를 $m:n$으로 내분하는 점 P의 좌표를 $x$라 하면,

\[
\overline{AP} = x - \fillblank{1.2cm}\,,\qquad
\overline{PB} = \fillblank{1.2cm} - x
\]
\[
(x-x_1):(x_2-x) = m:n \quad\Rightarrow\quad x = \fillblank{4.5cm}
\]
\end{sheetbox}

\section{연습 문제}

\begin{sheetbox}
\textbf{1.} 다음 두 점 사이의 거리를 구하시오.\\
\hspace*{1.6em}⑴ A($-3$), B($3$) \qquad ⑵ A($6$), B($-2$)
\writespace{2}

\vspace{8pt}
\textbf{2.} 두 점 A($-6$), B($8$)에 대하여 다음 점의 좌표를 구하시오.\\
\hspace*{1.6em}⑴ 선분 AB를 $2:5$로 내분하는 점 \qquad ⑵ 선분 AB의 중점
\writespace{3}
\end{sheetbox}

\section{사고의 가시화 --- See · Think · Wonder}

\noindent{\footnotesize\color{inksoft} 오늘 배운 `내분'을 떠올리며 아래 세 칸을 채워 보자.}\\[6pt]
\begin{tabularx}{\linewidth}{@{}XXX@{}}
\textbf{\footnotesize\color{indigoc} See --- 무엇을 보았나?} &
\textbf{\footnotesize\color{indigoc} Think --- 무엇을 생각했나?} &
\textbf{\footnotesize\color{indigoc} Wonder --- 무엇이 궁금한가?} \\[4pt]
\end{tabularx}
\noindent
\begin{minipage}[t]{0.32\linewidth}\writespace{4}\end{minipage}\hfill
\begin{minipage}[t]{0.32\linewidth}\writespace{4}\end{minipage}\hfill
\begin{minipage}[t]{0.32\linewidth}\writespace{4}\end{minipage}

\section{마무리 성찰}

\begin{sheetbox}
\begin{minipage}[t]{0.48\linewidth}
{\footnotesize\color{inksoft} 오늘 배운 것을 한 문장으로}
\writespace{2}
\end{minipage}\hfill
\begin{minipage}[t]{0.48\linewidth}
{\footnotesize\color{inksoft} 감사 일기 / 유념 한 줄}
\writespace{2}
\end{minipage}
\end{sheetbox}

\end{document}
